\documentclass[12pt]{article}
\usepackage[utf8]{inputenc}
\usepackage[spanish]{babel}
\usepackage{amsmath}
\usepackage{amsfonts}
\usepackage{amssymb}
\usepackage{geometry}
\geometry{letterpaper, margin=2.5cm}

\begin{document}

\begin{center}
    \Large \textbf{Evaluación de Examen 1: Unidad I} \\
    \large Física para Ciencias de la Salud (005-1713) \\
    \vspace{0.5cm}
    \textbf{Estudiante:} Luisielis Marin \quad \textbf{C.I.:} 33.627.901
    \vspace{0.3cm} \\
    \large \textbf{Calificación Final: 7/10 puntos}
\end{center}

\vspace{0.5cm}

\section*{Análisis de Resultados}

\subsection*{1. Ángulo plano (Conversión a radianes)}
\textbf{Procedimiento y resultado:} Plantea correctamente la conversión usando el factor $\frac{\pi}{180}$. Simplifica la fracción de forma directa dividiendo entre 15, obteniendo el valor exacto de $\frac{5\pi}{12}$ rad. También calcula la aproximación decimal, demostrando un entendimiento completo. \\
\textbf{Veredicto:} \textbf{Correcto} (2/2 pts).

\subsection*{2. Sistemas de coordenadas (Longitud del hueso)}
\textbf{Procedimiento y resultado:} La estudiante plantea la fórmula de distancia, calcula las diferencias en los ejes ($6$ y $8$), sus cuadrados ($36$ y $64$) y la suma de estos ($100$). Sin embargo, el procedimiento queda incompleto al \textbf{omitir el paso final de aplicar la raíz cuadrada} para obtener la longitud de $10$ cm. \\
\textbf{Veredicto:} \textbf{Incompleto} (1/2 pts).

\subsection*{3. Vectores (Fuerza resultante)}
\textbf{Procedimiento y resultado:} Aunque el planteamiento inicial es algo desordenado, la estudiante agrupa y suma correctamente las componentes $\hat{i}$ y $\hat{j}$, llegando al vector resultante exacto de $(30\hat{i} + 60\hat{j})$ N. \\
\textbf{Veredicto:} \textbf{Correcto} (2/2 pts).

\subsection*{4. Potencias de diez (Notación científica y prefijo)}
\textbf{Procedimiento y resultado:} Se comete un error matemático al expresar el número en notación científica, escribiendo $1,25 \times 10^{-6}$ m, lo cual no corresponde al valor decimal original. Adicionalmente, se omitió por completo la segunda parte de la pregunta, al no identificar ni nombrar el prefijo del S.I. \\
\textbf{Veredicto:} \textbf{Incorrecto} (0/2 pts).

\subsection*{5. Sistemas de unidades (Conversión de densidad)}
\textbf{Procedimiento y resultado:} Excelente resolución. La estudiante no solo demuestra conocer el factor de conversión global ($1000$), sino que además plantea de forma impecable los factores de conversión en cadena para validar su resultado final de $1030 \text{ kg/m}^3$. \\
\textbf{Veredicto:} \textbf{Correcto} (2/2 pts).

\vspace{0.5cm}
\section*{Comentarios Generales y Sugerencias}
La estudiante demuestra un buen manejo de los factores de conversión y del álgebra vectorial.
\begin{itemize}
    \item Se recomienda completar siempre los procedimientos hasta el final, asegurándose de aplicar todas las operaciones necesarias (como la raíz cuadrada en la Pregunta 2).
    \item Es importante repasar las reglas de notación científica para evitar errores al contar los decimales (Pregunta 4).
\end{itemize}

\end{document}